% Section 2
\section{Setting, targets and assumptions}\label{sec:setting}

\subsection{Data and the accounting identity}

There are units $i=1,\dots,n$. Unit $i$ has size $N_i$ (its number of
members), row shares $X_i\in\Delta^{R-1}$ and column shares
$T_i\in\Delta^{C-1}$, where $\Delta^{k-1}=\{x\in\R^k:x\ge0,\ \one^\top x=1\}$
is the probability simplex and $\one$ is the vector of ones. Rows are groups
(or choices in an earlier election), indexed by $r,s=1,\dots,R$; columns are
outcomes, indexed by $c=1,\dots,C$. When units are nested in districts,
$D_i\in\{1,\dots,D^\star\}$ is the district of unit $i$, and $V_i=(X_i,N_i,D_i)$
collects the design variables (without districts, $V_i=(X_i,N_i)$).

The latent rate $\beta_{irc}$ is the share of the row-$r$ members of unit $i$ who
fall in column $c$. Each row of the matrix $\beta_i=[\beta_{irc}]$ lies in the
simplex, so $\beta_i\in\Bset:=(\Delta^{C-1})^R$. The accounting identity is
\begin{equation}\label{eq:acc}
T_{ic}=\sum_{r=1}^R X_{ir}\beta_{irc},\qquad\text{that is,}\qquad T_i=\beta_i^\top X_i .
\end{equation}
For given margins $(x,t)$, the \emph{unit Duncan--Davis polytope}
\[
S(x,t)=\{b\in\Bset:\ b^\top x=t\}
\]
is the set of rate matrices consistent with one unit's margins, and
$h_S(u;x,t)=\max_{b\in S(x,t)}\langle u,b\rangle$ is its support function, where
$\langle u,b\rangle=\sum_{rc}u_{rc}b_{rc}$. With cell masses $n_{rc}=x_rb_{rc}$
the polytope is a transportation polytope, and its single-cell extremes are the
Fréchet--Hoeffding bounds $\max\{0,x_r+t_c-1\}\le n_{rc}\le\min\{x_r,t_c\}$.

\subsection{Three targets}\label{sec:targets}

Three related targets appear. Each result says which one it concerns. Write
$w_{ir}=N_iX_{ir}/\sum_jN_jX_{jr}$ for the share of row-$r$ members who live in
unit $i$, and $m(v)=\E[\beta\mid V=v]$ for the conditional mean of the rates.
\begin{itemize}
\item The \emph{actual} rate $B_{rc}=\sum_iw_{ir}\beta_{irc}$: what the row-$r$
members of this electorate did. It is JKST's target.
\item The \emph{model-implied} rate $\Bt_{rc}=\sum_iw_{ir}m_{rc}(V_i)$: what the
systematic pattern implies they did, at the observed design. It is the
conditional mean of $B$ given the design.
\item The \emph{superpopulation} rate $B^*_{rc}=\E[NX_r\beta_{rc}]/\E[NX_r]$,
used in the population identification results.
\end{itemize}
The actual and model-implied rates differ by a residual term,
$B_{rc}=\Bt_{rc}+E_{rc}$ with $E_{rc}=\sum_iw_{ir}e_{irc}$ and
$e_i=\beta_i-m(V_i)$. A linear target is $\varphi=\sum_{rc}a_{rc}B_{rc}$ for
fixed weights $a_{rc}$; loyalty, defection, demobilisation and the gap between
two groups' rates are examples. All inference in this paper is conditional on
the design and covers $B$ or $\Bt$; inference for $B^*$ would need the
sampling variation of the design and is not pursued.

\begin{reading}
$B$ is what these voters actually did. $\Bt$ is what the systematic pattern
implies they did; it ignores the local luck of each unit. A confidence interval
for $B$ must allow for that luck, so it is wider.
\end{reading}

\subsection{Linear contextual effects in homogeneous form}

Because $\one^\top x=1$ on the simplex, an affine function of $x$ can be written
as a homogeneous linear one: $a+b^\top x=(a\one+b)^\top x$. The pooled LCE model,
called M0, says that for each column $c$ there is a matrix $H_c\in\R^{R\times R}$
with
\begin{equation}\label{eq:lce}
\E[\beta_{irc}\mid X_i,N_i]=e_r^\top H_cX_i\qquad\text{on }\{X_i\in A\},
\end{equation}
where $e_r$ is the $r$-th unit vector and $A\subseteq\Delta^{R-1}$ is the compact
convex region on which the model is assumed. The default region is the
population hull $A=\conv\supp P_X$. Write $H=(H_1,\dots,H_C)$ and
$M_H(x)=[e_r^\top H_cx]_{rc}$ for the implied mean rate matrix at composition
$x$. Split each matrix into its symmetric and skew parts,
\[
H_c=Q_c+K_c,\qquad Q_c=\tfrac12(H_c+H_c^\top),\qquad K_c=\tfrac12(H_c-H_c^\top).
\]
Taking conditional expectations of \eqref{eq:acc}, the regression function is
\[
\tau_c(x):=\E[T_c\mid X=x]=x^\top H_cx=x^\top Q_cx ,
\]
because $x^\top K_cx=0$ for a skew matrix. The regression therefore reveals $Q$
and never sees $K$. The unknown $K=(K_1,\dots,K_C)$ lives in
\[
\Kspace=\Big\{(K_1,\dots,K_C):\ K_c^\top=-K_c,\ \textstyle\sum_cK_c=0\Big\},
\qquad\dim\Kspace=(C-1)\,R(R-1)/2 ,
\]
which is one number in $2\times2$, two in $2\times3$ and six in $3\times3$. We
write $k$ for the free coordinates $k_{c,rs}$ ($c<C$, $r<s$) of $K$, and, for
$R=2$, $\kappa_c=K_{c,12}$.

\paragraph{What $H$, $Q$ and $K$ mean.} $H_{c,rs}$ is the rate at which row-$r$
members choose column $c$ in a hypothetical unit made up entirely of row-$s$
members. $Q_{c,rs}$ is the average of ``$r$-members in pure-$s$ units'' and
``$s$-members in pure-$r$ units''; the data identify it. $K_{c,rs}$ is half the
gap between the two; the data never identify it. This is the R$\times$C version
of JKST's question: is turnout low in a mixed area because group A turns out
less, or because group B turns out less where A is numerous?

\paragraph{Moments.} The superpopulation weights are $\mu_r=\E[NX_r]>0$ and
$\mu_{rs}=\E[NX_rX_s]$. In finite-design statements the same symbols denote the
sample analogues with weights $p_i=N_i/\sum_jN_j$ (Section~\ref{sec:signs}).

\subsection{The model and three identified sets}

Let $P$ be the law of the observables $W=(N,X,T)$ (with $D$ when there are
districts), and $\mathbb Q$ a law of $(W,\beta)$. The model $\mathcal M(H;A)$ is the
set of laws $\mathbb Q$ such that
\begin{itemize}
\item[(S1)] $\beta\in\Bset$ and $T=\beta^\top X$ almost surely (the accounting identity);
\item[(S2)] $\E[\beta\mid V]=M_H(X)$ almost surely on $\{X\in A\}$ (linear contextual
effects, with no dependence on size);
\item[(S3)] $M_H(x)\in\Bset$ for every $x\in A$ (the linear model gives valid
probabilities on all of $A$).
\end{itemize}
Let $Q(P)$ be the unique symmetric matrices with $\E_P[T_c\mid X=x]=x^\top Q_cx$
(Lemma~\ref{lem:decomp} gives existence and uniqueness under the model). Three
sets of values of $K$ are used.
\begin{itemize}
\item The \emph{sharp} identified set
$\ThI(P;A)=\{K\in\Kspace:\ \exists\,\mathbb Q\in\mathcal M(Q(P)+K;A)$ whose
$W$-marginal is $P\}$. It uses the whole law of the data and conditions on
$(X,N)$.
\item $\ThX(P;A)$: the same with (S2) replaced by the weaker
$\E[\beta\mid X]=M_H(X)$. This is the set of JKST's Proposition~1.
\item The \emph{outer} set
$\ThO(P;A)=\{K\in\Kspace:\ (Q_c(P)+K_c)x\ge0$ elementwise for all $c$ and all
$x\in A\}$. It uses only the regression function and the requirement that the
implied rates be probabilities.
\end{itemize}
Proposition~\ref{prop:nesting} shows $\ThI\subseteq\ThX\subseteq\ThO$. Three
hulls also appear: the population hull $\conv\supp P_X$ (the default $A$ in
identification), the hull of the empirical law of the design (exact for
finite-design identification), and the sample hull used as an estimate of $A$.
The sample hull lies inside the population hull, so estimating it can only widen
the outer bound (Proposition~\ref{prop:consistency}(c)). Two conventions hold
throughout. First, $\ThO$ is called sharp only \emph{for the regression function}
(Proposition~\ref{prop:sharpreg}). Second, every ``sharp'' statement is relative
to continuous (real-valued) latent rates, not integer cell counts.

\subsection{Assumptions and what each implies}\label{sec:assumptions}

Each assumption is stated and then read in plain terms: what it rules in or
out, when it is credible, and what goes wrong if it fails.
Table~\ref{tab:assumption-use} shows which rungs use which assumption.

\begin{assumption}[A1: closed population]\label{ass:A1}
$T_i=\beta_i^\top X_i$ exactly, with the same $N_i$ people in both margins.
\end{assumption}
\begin{reading}
Both margins describe the same people: the same electorate in both elections,
or the same eligible population for group and outcome. A1 rules out turnover,
moving, deaths, new registrants and boundary changes. It is credible when both
margins come from one register and the categories include abstention and,
where relevant, ``entered'' or ``left''. If it fails, the Duncan--Davis sets and
the regression identity both shift, and every bound can exclude the truth.
Bounded turnover changes both the unit polytopes and the regression identity,
so it cannot be handled by widening the final interval; it is an open
extension. Every rung, including Duncan--Davis, needs A1.
\end{reading}

\begin{assumption}[A2: linear contextual effects on $A$]\label{ass:A2}
Each group's mean rate is linear in the unit's composition on $A$:
equation~\eqref{eq:lce}.
\end{assumption}
\begin{reading}
A group's behaviour may depend on its surroundings, but only along a straight
line: for example, loyalty rises steadily as the in-party share of the unit
rises. A2 rules out thresholds and tipping (a machine that mobilises only above
50\%) and U-shapes (competitive units mobilising both sides). It is the
\emph{only} identifying restriction of the linear rungs. If it fails, the
quadratic regression is misspecified and the bound can miss the truth by an
amount that does not shrink with more data. Some violations are invisible:
nonlinear departures that cancel in the aggregate leave the data unchanged while
moving the target (Proposition~\ref{prop:undetectable}). A cubic term or an empty
bound can reveal some failures; goodness of fit can never validate A2.
\end{reading}

\begin{assumption}[A3: mean independence from size]\label{ass:A3}
$\E[\beta\mid X,N]=\E[\beta\mid X]$, with $D$ added to both conditioning sets
under the district rungs.
\end{assumption}
\begin{reading}
Large and small units with the same mix behave alike. This matters because the
target weights units by size. If large urban units differ in loyalty from small
rural ones with the same mix, the size-weighted target is biased by
$\Cov(NX_r,\beta_{rc}-m_{rc}(X))/\E[NX_r]$. District baselines help when size
differs by district (Proposition~\ref{prop:D4}). A3 is partly testable through
interactions of $X$ with $\log N$ in the regression.
\end{reading}

\begin{assumption}[A4: support and extrapolation]\label{ass:A4}
$\supp(X)\subseteq A$, $A$ is compact, and A2 holds on all of $A$.
\end{assumption}
\begin{reading}
The straight lines are trusted over the whole region $A$, whose corners (the
most extreme units) set the constraints. A larger $A$ gives tighter bounds but
assumes more: taking the whole simplex assumes that the lines extend to pure
units that may not exist (Proposition~\ref{prop:simplex}). Outlying units at the
corners of the hull can make constraints bind spuriously; trimming the hull is
the check. Probability validity on $A$, condition (S3), is automatic for the
population hull (Lemma~\ref{lem:outer}(c)).
\end{reading}

\begin{assumption}[A5: rank]\label{ass:A5}
For every symmetric $R\times R$ matrix $\Psi$, $x^\top\Psi x=0$ for all
$x\in\supp_AP_X$ implies $\Psi=0$. Here $\supp_AP_X$ is the support of the
restriction of $P_X$ to $A$: the smallest closed set $F$ with $P(X\in A\setminus F)=0$.
\end{assumption}
\begin{reading}
The compositions vary enough to fit the quadratic regression. For two groups it
means at least three distinct compositions. For three groups it needs six
compositions that do not lie on a common conic; any five do, so five are never
enough (Remark~\ref{rem:A5}). A5 fails, or nearly fails, when shares barely vary.
Then $Q$ is weakly identified and the bound is wide and unstable. The condition
number of the design matrix is the check. The restricted support matters only
when $A$ carries less than full mass: for $X$ uniform on $[0,1]$ and
$A=\{1/2\}$, $\supp P_X\cap A=\{1/2\}$ but $\supp_AP_X=\emptyset$.
\end{reading}

\begin{assumption}[A6: source of randomness and independence]\label{ass:A6}
Given the design $\{(N_i,X_i,D_i)\}$ and the district effects, the unit rates
$\beta_i$ are random with mean $m(V_i)$, and the residuals $e_i=\beta_i-m(V_i)$ are
independent across units.
\end{assumption}
\begin{reading}
Every unit is observed, so the uncertainty does not come from sampling units. It
is model-based: given its characteristics, each unit's behaviour has a random
part. A shock common to a district is a fixed feature of this electorate: part
of the mean under the district rungs, a misspecification under the pooled one.
Independence is a substantive assumption. Neighbourhood spillovers within a
district violate it, and then the variances and residual allowances are too
small. Observing every unit does not remove dependence. The pooling test of
Section~\ref{sec:inference} tests mean shifts, not dependence.
\end{reading}

\begin{assumption}[A7: regularity for the normal limit]\label{ass:A7}
A unique active constraint at the optimum of the LP, $\mu_{rs}>0$, and $\ThO$ not
a singleton.
\end{assumption}
\begin{reading}
This is JKST's tie-breaking condition. With it the endpoint of the bound is a
smooth function of the regression coefficients and the delta method applies.
In R$\times$C, and especially with many districts, near-ties are the norm, so the
inference of Section~\ref{sec:inference} does not need it. If it fails, the endpoint
is only directionally differentiable, the delta method loses its justification
and the standard bootstrap is inconsistent (Theorem~\ref{thm:limit}). This is a
failure of the inference, not a bias in identification.
\end{reading}

\begin{assumption}[A8: shared contextual effects across districts]\label{ass:A8}
The contextual effects are common to all districts; only the intercepts vary.
\end{assumption}
\begin{reading}
The way composition shifts behaviour is the same in every district, while each
district may be more or less loyal across the board (parallel lines). A8
produces the intersection gain of the shared-effects rung. It fails if, say, a
strong incumbent changes the contextual gradient in one district, or if floor
and ceiling effects make an additive shift impossible near 0 or 1, which is most
likely in the lopsided districts that tighten the bound most. The bound can then
exclude the truth and still look narrow. The symmetric part of the restriction
is testable (Lemma~\ref{lem:D1}(b)); the shared $K$ is not.
\end{reading}

\begin{assumption}[A9: within-district richness]\label{ass:A9}
Under district-specific effects, A5 within each district. Under shared effects,
each district's support contains $R$ affinely independent compositions and the
stacked design has full column rank (Lemma~\ref{lem:D1}(d)).
\end{assumption}
\begin{reading}
Districts with few, similar units identify almost nothing. Under shared effects
the bar is lower: a district's baseline needs only $R$ affinely independent
compositions. Weak identification per district, combined with a maximum or
minimum over districts, biases plug-in bounds inward.
\end{reading}

\begin{assumption}[A10: asymptotic regime]\label{ass:A10}
One of: (R1) $D^\star$ fixed and the number of units per district $n_d\to\infty$;
(R2) $D^\star\to\infty$ with $n_d$ bounded; (R3) both grow.
\end{assumption}
\begin{reading}
Which limit theory applies depends on the design. It is a modelling choice, not
a fact. With ten polling places in each of ninety ridings (R2), the district
intercepts are incidental parameters and the plug-in intersection is badly too
narrow. Table~\ref{tab:regimes} in Section~\ref{sec:inference} says which result
holds in which regime.
\end{reading}

\paragraph{Assumptions not made.} No distributional form for $\beta$ (unlike
King's truncated normal), no independence between $\beta$ and $X$ (contextual
effects are the point), no stationarity over time, and no point-identifying
nonlinearity.

\begin{table}[t]
\centering\small
\begin{tabularx}{\linewidth}{@{}L{0.33\linewidth}Y@{}}
\toprule
Rung & Assumptions used \\
\midrule
Duncan--Davis (DD) & A1 only. \\
Sign rows on realised tables (N1) & A1; continuous unit tables; no regression, no sampling inference. \\
Pooled LCE (M0) & A1--A5; inference adds A6, and A7 only for the normal limit. \\
District rungs (M2, M1($\delta$), M1, L4, L5, L6) & A1--A4 within districts, A9 in place of A5, A8 for M1 and the rungs below it; inference adds A6 and the regime A10. \\
Monotone means (N2) & A1, A3, and the finite-support convention of Section~\ref{sec:N2}. \\
Sign rows under a linear rung & Those of the rung they are added to. \\
\bottomrule
\end{tabularx}
\caption{Which rungs use which assumption. The rungs are defined in Sections~\ref{sec:districts} and~\ref{sec:signs}.}\label{tab:assumption-use}
\end{table}
